\documentclass[a4paper,12pt,numbers=enddot]{scrreprt}	% numbers=enddot fa seguire il puntino dopo il numero di cap. noenddot lo toglie 
\usepackage[utf8]{inputenc}	%imposta la codifica di input; invece di "latin1", anche "utf8" va bene
\usepackage[small]{caption}	%per i cara piccoli
\usepackage[english]{babel}	%l?ultima lingua � quella predefinita
\usepackage{listings}		%codice sorgente
\lstset{frame=shadowbox, rulesepcolor=\color{blue}, breaklines=true, basicstyle=\small\ttfamily, columns=fullflexible,  keepspaces=true}	%disegna un rettangolo attorno al codice sorgente
\usepackage{indentfirst}	%rientra il primo paragrafo di ogni sezione
\usepackage{amsmath}%,amssymb,amsthm       % indispensabili per la matematica
\usepackage{graphicx}		% per inserire le immagini
\usepackage{graphics}		%?
\usepackage{sidecap}		%per mettere caption a lato della figuta
\usepackage{caption}		% didascalie personalizzate
\usepackage[numbers]{natbib} %?
\usepackage{psfrag} %?
\usepackage{multirow}	%?
\usepackage{booktabs} 	%?
\usepackage{lscape}%?
\usepackage{longtable}%?
\usepackage{enumitem}%?
\usepackage{epsfig}		%per mettere fig usando comando tabular
\usepackage{makeidx}	%indice
\usepackage{varioref}
\usepackage[table]{xcolor}
\usepackage{epstopdf}
\usepackage{wrapfig}
%Setta dimensioni body
\usepackage[total={16.5cm,25cm},
						top=3.5cm,
						bottom=3.5cm,
						left=2.5cm,
						right=2.5cm,
						%includefoot,
						centering]{geometry}

\usepackage{wallpaper}%?
\usepackage{color}%?
\usepackage{setspace}		%per definire con dopo l'interlinea con \onehalfspacing
\onehalfspacing
\usepackage{appendix}
\usepackage{acronym}%?
\usepackage{makecell}

%Intestazioni e pie? pag.  
\usepackage[automark,headsepline,footsepline,autooneside,plainheadsepline,plainfootsepline]{scrpage2}	%i comandi plain-ecccc servono a dire che nella prima pag. del nuovo cap. voglio queste cose.
\usepackage[pagebackref=true, colorlinks=true, linkcolor=blue, hyperindex=true, urlcolor=blue, citecolor = green, linktocpage=true]{hyperref}		% indici e riferimenti cliccabili
\pagestyle{scrheadings}
\clearscrheadfoot
\ihead [Enrico Rizzardi]{Enrico Rizzardi}	%Kopfzeile Innen
\chead {}					%Kopfzeile Mitte
\ohead [\headmark]{\headmark}			%Kopfzeile Au�en
\ifoot {}					%Fu�zeile Innen
\cfoot[\pagemark]{\pagemark} 			%Fu�zeile Mitte
\ofoot {}

\automark[chapter]{chapter}
\KOMAoptions{cleardoublepage=scrheadings}	%per lasciate style=scrheadings anche per le pag vuoto (svuotate col comando cleardoublepage), si puo? anche scegliere cleardoublepage=plain o empty

\footskip=30pt	 %	abbassa il margine inf della pag
\setlength{\headheight}{2\baselineskip}

\newcommand\SectionFontStyle{\rmfamily}			% geh�rt zu den beiden folgenden setkomafont , keine ahnung was es macht =)
\setkomafont{chapter}{\huge\bfseries\SectionFontStyle}	% Chapter in der gleichen schriftart SectionFontStyle
\setkomafont{sectioning}{\bfseries\SectionFontStyle}
%Crea i link nel pdf (vedere come si fa a farli verdino e con i collegamenti solo sui numeri, non questi rettangoli cessi rossi	

\makeindex
\begin{document}

\renewcommand{\bibname}{References}	%to change the name: bibliography to references

%\pagenumbering{roman}
%\pagestyle{plain}
%******************************************************************
% Materiale iniziale
%******************************************************************

\begin{titlepage}   %Frontespizio
\begin{center}
{\Large UNIVERSIT\`A DEGLI STUDI DI BRESCIA}

\vspace{0.2cm}

{\Large FACOLT\`A DI INGEGNERIA}

\begin{figure*}[htbp] 
\begin{center}
 \begin{tabular}{c c}
\raisebox{-.5\height}{\includegraphics[width=0.20\textwidth]{images/00-loghi/logo_UniBS.eps}}
 & \raisebox{-0.57\height}{\includegraphics[width=0.20\textwidth]{images/00-loghi/freiburg.png}}
\end{tabular}
\end{center}
\end{figure*}

% \vspace{0.5cm}

CORSO DI LAUREA IN INGEGNERIA INFORMATICA \\

TESI DI LAUREA SPECIALISTICA \\

%\vspace{0.8cm}
\vspace{1.5cm}
{\normalsize \emph{\textbf{ANALISI DI FATTIBILIT\`A ED IMPLEMENTAZIONE}}} \\

\vspace{0.2cm}

{\normalsize \emph{\textbf{DI UN SISTEMA DI ROBOTICA COGNITIVA}}} \\

\vspace{0.2cm}

{\normalsize \emph{\textbf{PER COMPITI DI NAVIGAZIONE}}} \\

%{\normalsize \emph{\textbf{IN REGIME INTERMITTENTE}}} \\

\vspace{0.5cm}

{\small \emph{\textbf{(Analysis of preconditions and implementation of a \\ Cognitive Robotic System for navigation tasks)}}} \\

\vspace{1.5cm}

\begin{flushleft}

{\normalsize Relatore:}
\begin{quote}
{\large \textbf{Ch.mo Prof. Riccardo Cassinis}} \\
\vspace{0.2cm}
\end{quote}

{\normalsize Correlatore:}
\begin{quote}
{\large \textbf{Ch.mo Prof. Marco Ragni}} \\
\vspace{0.2cm}
\end{quote}

\end{flushleft}

% \vspace{0.5cm}

\begin{flushright}
Laureando: \\
\textbf{Stefano Bennati} \\
\vspace{0.05cm}
\textbf{Matricola 72363} \\
\vspace{1.5cm}
\end{flushright}

\large{Anno Accademico 2010-2011}

\end{center}
\end{titlepage}

\pagenumbering{arabic}
\tableofcontents
\listoffigures

\newpage
\thispagestyle{empty}
\mbox{}

\chapter{Introduction}
\section{Background}
\section{Problem description}
\section{Working Instruments}
\subsection{ACT-R}
\subsection{OpenCV}
	OpenCV, \emph{Open Source Computer Vision}, is a library for computer vision. It was originally developed by Intel, and now it is supported by Willow Garage.
	It is a cross-platform library and it is released under a BSD license, thus it is free and open source. It was developed in C and C++ and then interfaces for Java and Python were added. OpenCV was designed for computational efficiency and with a strong focus on real-time applications. The version 2.4 has more than 2500 algorithms. The library has been used in many applications as, for example, mine inspection and robotics %~\cite{OCVMain}.
	\begin{wrapfigure}{l}{0.2\textwidth}
  		\begin{center}
  	  	\includegraphics[width=0.2\textwidth]{images/ch_01/opencv_logo.png}
  		\end{center}
  		\caption{OpenCV Logo}
	\end{wrapfigure}		
	
	\subsubsection*{History}
	The OpenCV Project started in 1999 as an Intel Reasearch initiative aimed to improve CPU intensive applications as a part of projects including real-time ray tracing and 3D display walls. The early goals of the project were developing optimized code for basic vision infrastructure, spreading this infrastructure to developers, and making it portable and available for free without forcing the developers to create their applications for free.
	The first alpha version was released to the public in 2000, followed by five beta versions between 2001 and 2005, which lead to version 1.0 in 2006. In 2008, the technology incubator Willow Garage begun supporting the project and, in the same year, version 1.1  was released. 
	On October 2009, OpenCV 2.0 was released. It included many improvements, such as a better C++ interface, more programming patterns, new functions and multicore optimized implementations of old ones. Willow Garage releases official versions of OpenCV every six months.
	\subsubsection*{Main Features}
		\includegraphics[width=1\textwidth]{images/ch_01/opencv_overview.jpg}
	The figure above shows the main features of the library. First of all, OpenCV offers an easy way to manage image data types. It provides functions to load, copy, edit, convert and store images and a basic graphical user interface that lets the developers handle keyboard and mouse and display image and video. The library can work both on images and videos and lets manipulate images even with matrix and vector algebra routines, for example, eigenvalues and singular value decomposition. It supports the most common dynamic data structures and offers many different basic image processing functions: filtering, edge and corner detection, color conversion, sampling and interpolation, morphological operations, histograms and image pyramids. Beyond this, it integrates many functions for structural analysis of the image including connected components, contour processing, distance transform, template matching, Hough transform, polygonal approximation, line fitting, ellipse fitting and Delaunay triangulation as well as features for camera calibration like finding and tracking calibration patterns, calibration, fundamental matrix estimation, homography estimation and stereo correspondence. It also offers many functions for motion analysis like optical flow, motion segmentation and tracking and features for object recognition.

\newpage
%\chapter{Project}
\section{SCRUM}
	in questo capitolo descriviamo lo scrum
%\input{Chapters/introduzione}

\pagenumbering{roman}


\newpage
\clearpage
\appendix
%\input{Chapters/appendice}
\newpage
\clearpage



%******************************************************************
% Materiale Finale
%******************************************************************
%******************************************************************
%References (two methods in Latex, this (bibTex) with the databe is the best
%******************************************************************
\cleardoublepage
\phantomsection	%gives hyperref something to hold on to when making the link.
\addcontentsline{toc}{chapter}{References}
\bibliographystyle{alpha}				%scelgo come formattare la biblio (abbrv abbrevia i nomi propri. plain non li abbrevia. Ce ne sono un sacco).
\bibliography{BeginningEnding/database_references.bib} 	%database dove prendere le info  
\nocite{*}						%stampa TUTTO quello presente nella banca dati, altrimenti solo la roba citata nella tesi

%NB in kile bisogna aggiungere BibTex al QuickBuild per far compilare la bibliografia. Metterlo subito dopo Latex.


\newpage										
\clearpage


\end{document}
